% from ega4-21a-fr.tex
\oldpage[IV-4]{261}
fractionnaire inversible $\sh{O}_U.f^{-1}$;
en vertu de \sref{IV.21.1.6}, pour toute fonction méromorphe régulière $f$, on a la relation
\[
\label{IV.21.2.8.2}
  f\in\Gamma(X,\sh{O}_X(D)) \iff \operatorname{div}(f)\geq-D.
  \tag{21.2.8.2}
\]
En outre, il est clair que l'on a des isomorphismes canoniques \sref{IV.21.2.3}
\[
\label{IV.21.2.8.3}
  \left\{\begin{aligned}
    &\sh{O}_X(0)=\sh{O}_X,\quad \sh{O}_X(D+D')=\sh{O}_X(D)\sh{O}_X(D')\xrightarrow{\sim}\sh{O}_X(D)\otimes\sh{O}_X(D')\\
    &\sh{O}_X(nD)=(\sh{O}_X(D))^n\xrightarrow{\sim}(\sh{O}_X(D))^{\otimes n}
  \end{aligned}\right.
  \tag{21.2.8.3}
\]
pour tout entier $n\in\mathbf{Z}$, et pour deux diviseurs quelconques $D$, $D'$ sur $X$.
\end{env}

\begin{env}[21.2.9]
\label{IV.21.2.9}
Soit $\sh{J}$ un Idéal fractionnaire inversible sur $X$.
L'injection canonique $\sh{J}\hookrightarrow\sh{M}_X$ définit par tensorisation un homomorphisme de $\sh{O}_X$-Modules
\[
\label{IV.21.2.9.1}
  \sh{M}_X(\sh{J})=\sh{J}\otimes_{\sh{O}_X}\sh{M}_X\to\sh{M}_X\otimes_{\sh{O}_X}\sh{M}_X=\sh{M}_X
  \tag{21.2.9.1}
\]
qui est un \emph{isomorphisme} : en effet, si $U$ est un ouvert de $X$ tel que $\sh{J}|U=\sh{O}_U.f$, où $f\in\Gamma(U,\sh{M}_X^*)$, l'homomorphisme \sref{IV.21.2.9.1} restreint à $U$ n'est autre que l'isomorphisme qui à toute section $t$ de $\sh{M}_X(\sh{J})|U=\sh{M}_X|U$ au-dessus de $V\subset U$, fait correspondre la section $t(f|V)^{-1}$ du même faisceau.
Dans l'isomorphisme \sref{IV.21.2.9.1}, aux sections méromorphes régulières de $\sh{J}$ au-dessus de $X$ correspondent les fonctions méromorphes régulières sur $X$.

Considérons en particulier le cas où $\sh{J}=\sh{O}_X(D)$, où $D$ est un diviseur sur $X$;
on a alors un isomorphisme canonique
\[
\label{IV.21.2.9.2}
  \sh{M}_X(\sh{O}_X(D))\xrightarrow{\sim}\sh{M}_X
  \tag{21.2.9.2}
\]
et on désigne par $s_D$ la section méromorphe régulière de $\sh{O}_X(D)$ au-dessus de $X$ qui correspond par cet isomorphisme à la section $1$ de $\sh{M}_X$.
Si $U$ est un ouvert de $X$ tel que $\sh{O}_X(D)|U=\sh{O}_U.f^{-1}$, où $f\in\Gamma(U,\sh{M}_X^*)$, on a $s_D|U=1$ dans $\Gamma(U,\sh{M}_X)$;
comme alors on a $D|U=\operatorname{div}_U(f)$, on en déduit \sref{IV.21.1.4} que l'on a
\[
\label{IV.21.2.9.3}
  \operatorname{div}(s_D)=D.
  \tag{21.2.9.3}
\]
D'autre part, on déduit aussitôt des isomorphismes canoniques \sref{IV.21.2.8.3} les formules
\[
\label{IV.21.2.9.4}
  s_0=1,\quad s_{D+D'}=s_D\otimes s_{D'},\quad s_{nD}=s_D^{\otimes n}\quad(n\in\mathbf{Z}).
  \tag{21.2.9.4}
\]
\end{env}

\begin{env}[21.2.10]
\label{IV.21.2.10}
Considérons, entre deux couples $(\sh{L},s)$, $(\sh{L}',s')$, où $\sh{L}$ et $\sh{L}'$ sont deux $\sh{O}_X$-Modules inversibles, $s$ (resp. $s'$) une section méromorphe régulière de $\sh{L}$ (resp. $\sh{L}'$) au-dessus de $X$, la relation : « il existe un isomorphisme $u:\sh{L}\xrightarrow{\sim}\sh{L}'$ tel que $\overline{u}(s)=s'$ », où $\overline{u}:\Gamma(X,\sh{M}_X(\sh{L}))\xrightarrow{\sim}\Gamma(X,\sh{M}_X(\sh{L}'))$ est l'isomorphisme déduit de $u$ (on notera que l'isomorphisme $u$ vérifiant cette condition est alors déterminé de façon unique).
Il est clair que c'est une relation d'équivalence, et puisqu'il existe un ensemble de $\sh{O}_X$-Modules inversibles tel que tout $\sh{O}_X$-Module inversible soit isomorphe à un élément de cet ensemble (\hyperref[0.5.4.7]{$\mathbf{0}_{\mathrm I}$, 5.4.7}), on peut parler de l'ensemble $D(X)$ des classes d'équivalence des couples $(\sh{L},s)$ pour la relation précédente.
Pour tout $\sh{O}_X$-Module inversible $\sh{L}$ et toute section
